\documentclass[11pt,a4paper]{article}
\usepackage[T1]{fontenc}
\usepackage[utf8]{inputenc}
\usepackage{newtxtext,newtxmath}
\usepackage[a4paper,margin=33mm]{geometry}
\usepackage{microtype}
\usepackage{graphicx}
\usepackage{xcolor}
\usepackage{mdframed}
\let\openbox\relax
\usepackage{amsmath,amsthm,amssymb,mathtools}
\usepackage{needspace}
\usepackage[explicit]{titlesec}
\usepackage{titling}
\usepackage[
  pdfusetitle,
  hidelinks,
  bookmarksopen=true,
  bookmarksnumbered=true
]{hyperref}

\linespread{0.95}
\setlength{\parindent}{0pt}
\setlength{\parskip}{5.5pt}
\allowdisplaybreaks
\numberwithin{equation}{section}

\DeclareMathOperator{\Tr}{Tr}
\DeclareMathOperator{\Var}{Var}
\DeclareMathOperator{\Cov}{Cov}
\DeclareMathOperator{\diag}{diag}
\DeclareMathOperator{\Sym}{Sym}
\DeclareMathOperator{\KL}{KL}
\newcommand{\R}{\mathbf R}
\newcommand{\E}{\mathbf E}

\titleformat{\section}
{\normalfont\normalsize\bfseries\raggedright}
{\thesection}{0.75em}{\begin{minipage}[t]{0.9\textwidth}#1\end{minipage}}
[\vspace{0.12em}\titlerule]

\newtheoremstyle{thmcompact}
{0.5ex}{0.5ex}
{\itshape}
{0pt}
{\bfseries}
{.}
{0.6em}
{\thmname{#1}\ \thmnumber{#2}\ \thmnote{(\textit{#3})}}

\theoremstyle{thmcompact}
\newtheorem{theorem}{Theorem}[section]
\newtheorem{proposition}[theorem]{Proposition}
\newtheorem{corollary}[theorem]{Corollary}
\newtheorem{definition}[theorem]{Definition}
\newtheorem{lemma}[theorem]{Lemma}

\theoremstyle{remark}
\newtheorem*{remarkplain}{Remark}

\mdfdefinestyle{thmframe}{%
leftline=true, rightline=false, topline=false, bottomline=false,
linewidth=0.8pt, linecolor=black,
innerleftmargin=16pt, innerrightmargin=8pt,
innertopmargin=6pt, innerbottommargin=6pt,
skipabove=6pt, skipbelow=6pt,
nobreak=true
}
\surroundwithmdframed[style=thmframe]{theorem}
\surroundwithmdframed[style=thmframe]{proposition}
\surroundwithmdframed[style=thmframe]{corollary}
\surroundwithmdframed[style=thmframe]{definition}
\surroundwithmdframed[style=thmframe]{lemma}

\newmdenv[
linewidth=0.6pt,
topline=false,
bottomline=false,
rightline=false,
skipabove=\baselineskip,
skipbelow=\baselineskip,
backgroundcolor=gray!3,
leftmargin=0pt,
innerleftmargin=8pt,
innerrightmargin=6pt,
frametitleaboveskip=0.4em,
frametitlebelowskip=0.2em,
frametitlefont=\bfseries\small,
]{remarkbox}
\newenvironment{remark}[1][]%
{\begin{remarkbox}\begin{remarkplain}[#1]\normalfont}
{\end{remarkplain}\end{remarkbox}}

\renewcommand{\qedsymbol}{}
\newcommand{\thmneed}{\Needspace{6\baselineskip}}
\AtBeginEnvironment{theorem}{\thmneed}
\AtBeginEnvironment{proposition}{\thmneed}
\AtBeginEnvironment{corollary}{\thmneed}
\AtBeginEnvironment{definition}{\thmneed}
\AtBeginEnvironment{lemma}{\thmneed}

\title{\textbf{Weakpoints in the Branch and Non-Gaussian Layer}\\
\large{Fixed fibre precision, entropic volume, noncommuting Hessians, and verdict-changing cumulants}}
\author{J.\,R.~Dunkley}
\date{11 April 2026}

\hypersetup{
  pdftitle={Weakpoints in the Branch and Non-Gaussian Layer},
  pdfauthor={J. R. Dunkley},
  pdfsubject={Weakpoint research for branch selection and non-Gaussian observation geometry},
  pdfkeywords={observation geometry, affine-hidden sector, entropic quotient, branch selection, cumulants}
}

\begin{document}

\maketitle

\begin{abstract}
This addendum stress-tests the strongest claims in \emph{Branch Selection and Non-Gaussian Findings}. The result is not a rollback. It sharpens the exact boundaries. The affine-hidden sector is exact because fixed hidden precision removes the fibre-volume term from branch comparison. If the hidden fibre precision depends on the visible variable, the exact law acquires a visible $\frac12\log\det D(v)$ correction, which can change the branch verdict. The affine-hidden tower law survives arbitrary staged elimination by Schur associativity. The branch selector Hessian does not require commuting blocks. The variational quotient and entropic quotient separate at order $\varepsilon$ through the Laplace fibre-volume term. Finally, fibre cumulants give a minimal explicit full-law failure mode: a branch can be quadratically tied while the actual visible law already moves at first cumulant order.
\end{abstract}

\section{Purpose}

The findings note already gives the positive story: exact local quadratic geometry, exact Gaussian full-law closure, fixed-shape non-Gaussian survival, and the affine-hidden exact sector. This addendum records the weakpoints that matter if one tries to promote those results too far.

The guiding rule is:
\[
\text{exact quadratic branch result}
\not\Rightarrow
\text{exact arbitrary non-Gaussian branch law}.
\]
The interesting part is not that the implication fails. The interesting part is where it fails and which extensions remain exact after the failure is understood.

\section{Fixed hidden precision is a real hypothesis}

The affine-hidden sector in the findings note assumes
\[
p(v,h)\propto
\exp\!\left(-A(v)-\frac12h^{\mathsf T}Dh-J(v)^{\mathsf T}h\right),
\qquad D\succ0,
\]
with $D$ independent of $v$. This is exactly why marginalisation and minimisation agree up to a constant.

\begin{proposition}[Variable hidden precision boundary]
\label{prop:variable-precision}
Let
\[
p(v,h)\propto
\exp\!\left(-A(v)-\frac12h^{\mathsf T}D(v)h-J(v)^{\mathsf T}h\right),
\qquad D(v)\succ0.
\]
Then the exact visible negative log-density is
\begin{equation}
S_{\mathrm{vis}}(v)
=A(v)+\frac12\log\det D(v)
-\frac12J(v)^{\mathsf T}D(v)^{-1}J(v)+\mathrm{const}.
\label{eq:variable-precision-visible}
\end{equation}
The variational hidden elimination gives
\begin{equation}
S_{\mathrm{var}}(v)
=A(v)-\frac12J(v)^{\mathsf T}D(v)^{-1}J(v).
\label{eq:variable-precision-var}
\end{equation}
Thus the exact marginal and the variational quotient differ by the visible fibre-volume term
\[
\frac12\log\det D(v).
\]
\end{proposition}

\begin{proof}
Complete the square in $h$ at each fixed $v$:
\[
\frac12h^{\mathsf T}D(v)h+J(v)^{\mathsf T}h
=
\frac12(h+D(v)^{-1}J(v))^{\mathsf T}D(v)(h+D(v)^{-1}J(v))
-\frac12J(v)^{\mathsf T}D(v)^{-1}J(v).
\]
The Gaussian integral over the fibre now contributes $(\det D(v))^{-1/2}$, which is not constant in $v$.
\end{proof}

\begin{corollary}[Minimal hostile example]
\label{cor:hostile-variable-precision}
Set $A(v)=0$, $J(v)=0$, and $D(v)=e^{2v}$ on an interval. Then $S_{\mathrm{var}}$ is flat, but
\[
S_{\mathrm{vis}}(v)=v+\mathrm{const}.
\]
The exact law selects the smallest allowed $v$ while the variational quotient cannot distinguish any $v$.
\end{corollary}

\begin{remark}
This is the cleanest weakpoint found in the affine-hidden story. The larger variable-precision sector is still exactly integrable, but it is not the same minimisation/marginalisation coincidence theorem. If it becomes part of the theory, it should be named as an entropic fibre-volume sector, not folded silently into the fixed-precision affine-hidden sector.
\end{remark}

\section{Affine-hidden tower law is stronger than two blocks}

The tower law is not just a two-block trick. It is Schur-complement associativity.

\begin{proposition}[Permutation-invariant staged hidden elimination]
\label{prop:permutation-tower}
Let $D\succ0$ be a fixed hidden precision on $\R^r$ and let $J(v)\in\R^r$. For any ordered partition of the hidden coordinates into blocks, staged completion of the square gives the same visible action
\[
A(v)-\frac12J(v)^{\mathsf T}D^{-1}J(v)+\mathrm{const}
\]
as eliminating all hidden coordinates in one step.
\end{proposition}

\begin{proof}
Each staged elimination is a block Gaussian integration and replaces the remaining hidden precision by the Schur complement of the eliminated block. Schur complements of a positive definite quadratic form are associative under ordered block elimination. The coupling vector transforms by the same block solve, and the accumulated constant term is exactly $-\frac12J^{\mathsf T}D^{-1}J$.
\end{proof}

The numerical stress test eliminated five hidden coordinates in all $5!$ orders across a nonlinear grid of $J(v)$ values. The maximum residual against one-step elimination was $2.78\cdot10^{-16}$.

\section{Branch Hessian does not need commuting blocks}

The selector Hessian theorem only needs exact block closure at the branch. It does not require the blocks $\{A_a\}$ or $\{B_a\}$ to commute.

\begin{proposition}[Noncommuting block stress]
\label{prop:noncommuting-hessian-stress}
For an exact closure branch $U\oplus W$ with block-diagonal symmetric family
\[
D_a=
\begin{pmatrix}
A_a&0\\
0&B_a
\end{pmatrix},
\]
the second variation remains
\[
\delta^2F_{\mu,U}[X]
=2\langle X,(M_WX-XM_U)\rangle
-2(1+\mu)\sum_a\|B_aX-XA_a\|_F^2,
\]
even when $[A_a,A_b]\ne0$ and $[B_a,B_b]\ne0$.
\end{proposition}

\begin{proof}
The derivation uses only the graph expansion of the projector, the block-diagonal form, and the identity
\[
\Tr(\Pi M)=\mathcal S(\Pi)+\mathcal L(\Pi),
\qquad M=\sum_aD_a^2.
\]
No simultaneous diagonalisation is used.
\end{proof}

The numerical stress test used two noncommuting $2\times2$ blocks on both sides. The internal commutator norms were $1.97$ and $2.40$, while the finite-difference Hessian agreed with the analytic Hessian to relative Frobenius error $1.52\cdot10^{-6}$.

\section{Variational and entropic quotients split at finite noise}

Let
\[
S_{\mathrm{ent},\varepsilon}(v)
:=-\varepsilon\log\int \exp(-S(v,h)/\varepsilon)\,dh.
\]
When the hidden minimum is nondegenerate, Laplace theory gives
\[
S_{\mathrm{ent},\varepsilon}(v)
=S_{\mathrm{var}}(v)
+\frac{\varepsilon}{2}\log\det S_{hh}(v,h_*(v))
+O(\varepsilon^2)
\]
up to constants independent of $v$.

The finite-noise stress test used a nonquadratic one-dimensional fibre. After centering by constants, the sup error of the Laplace-corrected approximation was:
\[
\begin{array}{c|c|c}
\varepsilon & \text{Laplace corrected error} & \text{variational-only error}\\
\hline
0.16 & 4.66\cdot10^{-4} & 9.71\cdot10^{-3}\\
0.08 & 1.23\cdot10^{-4} & 4.96\cdot10^{-3}\\
0.04 & 3.19\cdot10^{-5} & 2.51\cdot10^{-3}.
\end{array}
\]
The corrected error ratios were about $3.78$ and $3.88$ under halving $\varepsilon$, consistent with an $O(\varepsilon^2)$ remainder. The variational-only error is much larger and scales as the missing $O(\varepsilon)$ fibre-volume term.

\section{Cumulants can change a tied quadratic branch verdict}

The cubic hidden-visible pathology can be stated as a branch verdict failure.

\begin{proposition}[Quadratic tie, full-law separation]
\label{prop:cumulant-verdict}
Let the baseline be
\[
S_0(x,z)=\frac12(x^2+z^2)+\alpha(x^4+z^4)
\]
and perturb one candidate branch by $\varepsilon xz^2$ while leaving another candidate unperturbed. At the origin the quadratic Hessian data of the two candidates are tied. But if $m_2=\E_0Z^2$, the perturbed branch has
\[
\E_\varepsilon X=-m_2^2\varepsilon+O(\varepsilon^2),
\qquad
\KL_{\mathrm{sym}}(p_\varepsilon^X,p_0^X)
=m_2^3\varepsilon^2+O(\varepsilon^3).
\]
Thus the full visible law separates the branches while the quadratic branch datum reports no difference.
\end{proposition}

\begin{proof}
The visible log-ratio is the conditional cumulant expansion
\[
\log\frac{p_\varepsilon^X(x)}{p_0^X(x)}
=\log\E_0[e^{-\varepsilon xZ^2}\mid x]
-\log\E_0[e^{-\varepsilon XZ^2}].
\]
The first nonconstant term is $-\varepsilon m_2x$. This gives the mean shift and the symmetric KL coefficient by the standard small log-density-ratio expansion.
\end{proof}

Numerically, at $\varepsilon=0.18$ the visible symmetric KL was $0.0027934$ against the cumulant prediction $0.0027580$, and the mean shift was $-0.0350396$ against the first-order prediction $-0.0348313$.

\section{Survival example}

The same stress test also included a branch-selection example fully inside the fixed-precision affine-hidden sector. Two nonlinear visible actions and nonlinear couplings were reduced exactly by
\[
S_i(v)=A_i(v)-\frac12D^{-1}J_i(v)^2.
\]
The exact branch gap was $0.02430$ in favour of branch $1$. There is no hidden caveat here: because $D$ is fixed, the law-level branch comparison is the same as the variational comparison up to an irrelevant constant.

\section{Module implications}

The module should not add a generic non-Gaussian branch selector on the basis of these results.

The following additions would be theorem-local and safe:
\begin{enumerate}
\item a branch selector Hessian diagnostic for exact closure branches;
\item a fixed-precision affine-hidden reducer;
\item a variable-precision Gaussian-fibre reducer that explicitly includes $\frac12\log\det D(v)$;
\item an experimental Laplace quotient helper clearly labelled as asymptotic;
\item fibre-cumulant diagnostics as a higher-order research layer, not as a current exact full-law engine.
\end{enumerate}

The main public boundary should remain: exact on supplied quadratic/Hessian/Fisher objects; exact in Gaussian full-law mode; exact in separately proved law sectors such as fixed-shape tilts and fixed-precision affine-hidden fibres; otherwise higher-order or global law structure is additional state.

\section{Conclusion}

The weakpoints are now sharper, not broader. The biggest one is the fibre-volume term. Fixed hidden precision is what makes affine-hidden marginalisation and variational elimination coincide. Visible-dependent hidden precision is still exactly integrable, but it belongs to the entropic quotient side and can change branch selection.

The second weakpoint is law leakage through cumulants. A branch can be tied at the quadratic layer and separated at the visible law layer by a first fibre cumulant. This is the cleanest reason not to promote the quadratic selector to an arbitrary non-Gaussian branch selector.

The wins remain substantial: affine-hidden towering is exact by Schur associativity, branch Hessians do not require commuting blocks, and the Laplace correction gives a precise bridge between variational and entropic quotients.

\end{document}
